\subsection{NVOS: source-calibrated response decisions}
N02 asks whether foreground is already ranked well but missed by a fixed
decision threshold. All three arms use the same final coverage/high field from
N01, reproduced with the same 120 updates. Every fourth sorted eligible source
view is a calibration view. One estimator renders these views with the all-source
field; another uses a field fitted without their reconstruction constraints.
Both use identical teacher targets, confidence weights and their common supported
pixels. The latter field is used only to estimate a threshold, not for target
prediction. Geometry and teacher tracking are shared: this is a reconstruction
holdout, not independent teacher or scene validation.

For source response $r_j$, soft target $y_j$ and weight $w_j$, we maximize
\begin{equation}
\widehat J(t)=\frac{\sum_j w_j y_j [r_j\geq t]}
{\sum_j w_j y_j+\sum_j w_j(1-y_j)[r_j\geq t]}.
\end{equation}
This is a ratio of soft expected counts, not expected IoU. It assumes teacher
evidence is useful for estimating a decision, without claiming calibrated
teacher correctness. All queries use the fixed grid $t=k/512$, $k=1,\ldots,511$;
ties prefer the value nearest 0.5 and then smaller $k$. Zero positive mass falls
back to 0.5. Excluding zero keeps unsupported zero responses negative.

\begin{table}[ht]
\centering\small
\resizebox{\linewidth}{!}{\begin{tabular}{lrrrrr}
\toprule
Task & Fixed 0.5 & Fit-source & Held-source & $t_{\rm fit}$ & $t_{\rm hold}$ \\
\midrule
fern & 83.53 & 84.31 & 84.27 & 0.3516 & 0.3652 \\
flower & 87.38 & 87.80 & 87.61 & 0.4492 & 0.4727 \\
fortress & 97.46 & 97.49 & 97.51 & 0.4941 & 0.4902 \\
horns\_center & 95.55 & 95.62 & 95.63 & 0.3555 & 0.3613 \\
horns\_left & 94.43 & 94.65 & 94.65 & 0.3496 & 0.3477 \\
leaves & 95.06 & 94.53 & 94.54 & 0.5957 & 0.5938 \\
orchids & 70.05 & 71.08 & 71.10 & 0.4199 & 0.4180 \\
trex & 78.75 & 74.63 & 75.35 & 0.3867 & 0.4023 \\
\midrule Mean (retained GT) & 87.78 & 87.52 & 87.58 & -- & -- \\
\midrule Mean (LUDVIG GT) & 87.93 & 87.58 & 87.65 & -- & -- \\
\bottomrule
\end{tabular}
}
\caption{Same NVOS continuous field, three decision rules. Thresholds are
estimated from source evidence before any target labels are read. All eight
responses and 24 decisions are sealed jointly. Continuous-response AP is
unchanged across these arms; no per-task winning rule is selected.}
\label{tab:nvos-source-readout}
\end{table}

This deliberately changes N01's fixed-threshold contract only for the declared
calibrated arms. Other benchmark outputs remain unchanged references. The known
cohort remains development data, and source improvement alone is insufficient
to establish better target responses.

The fixed readout reproduces N01 at 87.7776 IoU. Fit-source calibration obtains
87.5151 and reconstruction-holdout calibration 87.5836. The latter improves six
tasks but harms leaves and trex, with a mean change of $-0.1940$ points
(descriptive seven-scene bootstrap interval $[-1.2740,0.5191]$). All source
criteria improve at their selected thresholds, yet target precision decreases
from 95.8301 to 94.8640 while recall rises from 91.4323 to 92.2555.
Using the alternate GT convention gives 87.9280, 87.5762 and 87.6462 respectively;
the held-source arm then improves four tasks and harms four.

At trex, threshold 0.40234 recovers 30,527 foreground pixels but introduces
118,640 additional false positives, decreasing IoU by 3.3978 points. At orchids,
threshold 0.41797 recovers 20,859 foreground pixels with 1,013 additional false
positives, increasing IoU by 1.0537 points. At leaves, threshold 0.59375 removes
648 false positives but loses 4,375 foreground pixels. These counts are at native
GT resolution. This source-estimated scalar decision does not distinguish the
beneficial and harmful spatial changes in these cases. The reconstruction holdout does
not resolve the source-to-target mismatch in this experiment. We retain the
fixed N01 readout; source-estimated thresholds are not promoted, and per-task
winners are not substituted into the method score.

\paragraph{Foreground recovery versus background ranking.}
A separate post-seal audit uses the exact same coverage responses and freshly
renders all eight v10 responses. Changing foreground and background score groups
separately, then symmetrizing their marginal AP effects, gives an exact
score-space decomposition of the AP change. The GT-dependent hybrid scores are
diagnostics only, never deployable outputs. For orchids, the foreground change
contributes +0.2212 AP points and the background change $-2.2373$, totaling
$-2.0161$. Background pixels in $[0.01,0.05)$ increase from 94,415 to 197,928,
while binary false positives change only from 1,222 to 1,265. Thus a check only
at threshold 0.5 misses substantial ranking contamination. Trex's AP gain is
primarily associated with foreground score changes (+2.0713 points).
The full-cohort decomposition and error images are in the supplement. These
observations motivate spatially selective correction that preserves background
rejection; they do not yet isolate which upstream teacher or geometry operation
causes the score shifts.
